% Cote 119: offprints, reports, courses, CV and letters (typed and printed part): archive cote 119, site page 248 (PDF page 249).
% typescript. Draft compared: none.
% AI-transcribed by Claude Opus 5.5 (claude-opus-5-5, adaptive thinking, effort high), Message Batches API, batch r119_full (msgbatch_01DwZNm5Y7YsKBWyAN4sFs8w), 2026-11-13; not yet checked by eye.
\scanpage{119/248}
\mspage{8}

bien faite'' par opposition à la ``tête bien pleine''. Cela ne suppose pas
nécessairement que le maître-enseignant soit au courant, à chaque moment,
de tous les plus récents progrès de sa spécialité, ce qui demanderait de sa
part un effort de mise au courant souvent incompatible avec ses tâches
d'enseignement. Mais les longues vacances universitaires lui permettent
tout au moins de se mettre au courant périodiquement des plus importants
parmi ces progrès, ou de contribuer, dans la mesure de ses moyens, à
l'avancement de sa science. Ces moyens personnels sont évidemment très
variables de l'un à l'autre. Au point de vue du maître-enseignant où nous
nous plaçons ici, l'étendue des dons pour la recherche devient d'ailleurs
un point secondaire, dans la mesure où sa participation à l'édification de
la science n'est pas conçue comme une fin en soi, mais comme un moyen parmi
d'autres pour garder vivant l'enseignement,- étant entendu que cela suppose
d'abord que l'enseignant lui-même reste intellectuellement vivant, c'est-à-%
dire intellectuellement actif.

Dans tout cela, la vocation principale reste l'enseignement, et
en définitive \uline{le sujet principal d'intérêt du maître-enseignant est} l'\uline{étudiant}
\uline{lui-même}. Il est clair qu'une telle vocation ne pourra s'épanouir pleinement
que sous des conditions qui permettent un contact vivant, constant et étroit
avec les étudiants. On sait combien nous en avons été loins tout au long de
ces dernières années, avec des amphithéâtres-forums où le maître-gardien-de-%
la-science, seul devant la foule des aspirants-à-la-science, semble sorti
tout droit d'une hallucinante vision de Kafka.


5. Les développements qui précèdent démontrent sans doute suffisamment la
nécessité de distinguer entre le maître qui enseigne et le maître qui édifie
la science, et mettent en évidence quelques uns de leurs traits caractéris-%
tiques. Il doit être bien entendu qu'il ne peut être question d'établir entre
ces deux espèces de vocations une hiérarchie quelconque. Il convient sim-%
plement de réagir enfin à d'anciennes habitudes de pensée héritées du passé,
qui minimisaient le rôle du premier en le subord\over{\unsure{o}}{o}\inserted{n}nant au second. Il faut
donc réhabiliter la vocation pédagogique dans l'Université, en la mettant
à pied d'égalité avec la vocation de recherche, et en lui fournissant les

\margin{248}
